\capitulo{4}{Técnicas y herramientas}

\section{Entorno de desarrollo web}

El proyecto se desarrolla como una aplicación web ejecutada en el navegador. Este enfoque permite que el editor pueda utilizarse sin instalar una aplicación de escritorio específica, manteniendo gran parte de la lógica dentro del propio cliente web.

El entorno de desarrollo se apoya en Node.js y npm para la gestión de dependencias y la ejecución de los scripts del proyecto. A través de estos scripts se realizan tareas habituales como arrancar la aplicación en modo desarrollo, generar la versión de producción o ejecutar las pruebas definidas.

Este planteamiento encaja con el tipo de herramienta desarrollada, ya que el usuario interactúa directamente con una interfaz gráfica y con un área de edición visual. Por ello, gran parte del trabajo se centra en coordinar la interfaz, la representación gráfica del diagrama y el modelo interno que mantiene la aplicación.

\section{Librerías principales del editor}

La interfaz del proyecto se desarrolla utilizando React, una librería de JavaScript orientada a la construcción de interfaces de usuario mediante componentes \cite{react}. En este trabajo, React sirve como base para organizar la aplicación, gestionar sus distintas partes visuales y coordinar la interacción del usuario con el editor.

Para la representación y manipulación de los diagramas se utiliza mxGraph \cite{mxgraph}. Esta librería proporciona la base gráfica necesaria para trabajar con vértices, aristas y eventos dentro del área de edición. En el contexto del proyecto, mxGraph no se utiliza únicamente para dibujar elementos, sino también como parte del flujo que relaciona la representación visual del diagrama con el modelo interno de la aplicación.

Además, el proyecto utiliza Material UI para algunos componentes de la interfaz \cite{material-ui}. Esta librería facilita la creación de elementos visuales comunes, como botones, diálogos o controles de interacción, manteniendo una apariencia más uniforme dentro de la aplicación.

En conjunto, estas librerías permiten construir un editor web en el que la interfaz, el área gráfica y las acciones del usuario trabajan de forma coordinada. En este trabajo, lo relevante no es solo el uso de estas herramientas por separado, sino la forma en la que se integran para representar diagramas Entidad-Relación de manera coherente.

\section{Herramientas de validación y pruebas}

El proyecto utiliza distintas herramientas para revisar el comportamiento de la aplicación y mantener una calidad mínima en el código. Estas herramientas no sustituyen a la revisión manual del editor, especialmente en los aspectos visuales o de interacción, pero ayudan a detectar errores y regresiones durante el desarrollo.

Para las pruebas unitarias se utiliza Vitest \cite{vitest}. Esta herramienta permite comprobar partes concretas de la lógica de la aplicación, como funciones de validación, transformación o generación de resultados, sin necesidad de ejecutar manualmente todos los flujos desde la interfaz.

Además, el proyecto cuenta con pruebas end-to-end mediante Playwright \cite{playwright}. Estas pruebas permiten ejecutar escenarios completos sobre la aplicación en el navegador, simulando acciones del usuario y comprobando que determinados flujos siguen funcionando correctamente.

Junto a las pruebas, se utiliza Biome como herramienta de formateo y análisis estático del código \cite{biome}. Su uso ayuda a mantener un estilo más uniforme y a detectar ciertos problemas antes de integrar cambios en el proyecto. También se emplea Lefthook para ejecutar comprobaciones asociadas al flujo de trabajo con Git.

\section{Control de versiones y apoyo al desarrollo}

Para el control de versiones se utiliza Git, lo que permite registrar los cambios realizados en el proyecto y trabajar de forma más ordenada durante el desarrollo \cite{git}. El uso de commits y ramas facilita separar las distintas tareas y mantener un historial del trabajo realizado.

El proyecto se gestiona mediante GitHub, que se utiliza como repositorio remoto y como apoyo para la organización del desarrollo \cite{github}. A través de sus herramientas se pueden documentar tareas, revisar cambios y mantener una trazabilidad básica de la evolución del editor.

Además de las herramientas propias del código, durante la redacción de la memoria se utiliza Overleaf como entorno de trabajo para el proyecto LaTeX. Esto permite editar la documentación de forma progresiva y comprobar que la memoria y los anexos compilan correctamente mientras se van completando los distintos capítulos.